problem:  
Let $\Sigma_{p,q}^7$ be a Milnor exotic 7‑sphere realized as the total space of an $S^3$‑bundle over $S^4$, corresponding to clutching data $(p,q)\in\pi_3(SO(4))\cong\mathbb{Z}\oplus\mathbb{Z}$ with $p+q=1$.  
Denote by $S(T\Sigma_{p,q}^7)$ the **sphere bundle of the tangent bundle** of $\Sigma_{p,q}^7$ (i.e. the bundle of unit tangent directions, canonically defined up to diffeomorphism).  
Let $g_{p,q}$ range over all Riemannian metrics on $\Sigma_{p,q}^7$, and let each induce a **connection metric** $\tilde g_{p,q}$ on $S(T\Sigma_{p,q}^7)$ via its Levi‑Civita connection, so that $\pi:(S(T\Sigma_{p,q}^7),\tilde g_{p,q})\to(\Sigma_{p,q}^7,g_{p,q})$ is a Riemannian submersion with totally geodesic $S^6$ fibers.  

**Problem:**  
Determine whether the topology of the clutching data $(p,q)$ can obstruct the existence of a connection metric $\tilde g_{p,q}$ with nonnegative sectional curvature on $S(T\Sigma_{p,q}^7)$.  
Equivalently, do there exist distinct Milnor spheres $\Sigma_{p,q}^7,\Sigma_{p',q'}^7$ such that their tangent sphere bundles are diffeomorphic, yet one admits a connection metric of nonnegative curvature while the other provably does not?  
As a refinement, can any curvature‑based obstruction—derived from characteristic classes, the Eells–Kuiper invariant, or O’Neill curvature tensors—distinguish the smooth class of $\Sigma_{p,q}^7$ via the existence or nonexistence of such nonnegatively curved connection metrics?

Why is it a "good" problem:  
This formulation isolates a sharply posed geometric‑topological obstruction problem within a concrete and computable class of manifolds. The Milnor 7‑spheres are the simplest known examples where smooth structure varies without changing homeomorphism type, and the tangent sphere bundle $S(T\Sigma_{p,q}^7)$ lies at the interface between the base’s smooth structure and submersion curvature geometry.  

The problem asks whether the **Levi‑Civita horizontal distribution**—which governs the O’Neill curvature terms—can reflect the exotic structure encoded in $(p,q)$. Because these curvature tensors couple the base’s Riemannian geometry to vertical fiber geometry, understanding how (and whether) exotic smoothness can obstruct nonnegative sectional curvature would reveal new mechanisms connecting **differential topology** (through the clutching map and invariants like the Eells–Kuiper $\mu$‑invariant) with **Riemannian submersion theory** and **bundles of nonnegative curvature** studied by Grove, Ziller, and others.  

The question is simple, explicit, and based on concrete, low‑dimensional geometry, yet any resolution would require fundamentally new insight into how curvature interacts with the fine structure of differentiable manifolds—precisely the kind of difficulty and cross‑field depth that characterizes a strong open research problem.